% Source: https://artofproblemsolving.com/community/c3h3297446_how_do_computers_calculate_irrational_exponents
% Stable identity: aops-topic-3297446-q01
The real exponential is defined by the convergent series $\exp(x)=\sum_{j=0}^{\infty}\dfrac{x^j}{j!}$. Thus one can approximate both the input $e=\sum_{k=0}^{\infty}\dfrac{1}{k!}$ and the outer exponential by finite sums with rigorous tails.
For $0\le x<3$ and $N\ge2$, successive omitted terms have ratio at most $\dfrac{3}{N+2}$. Consequently
\[0\le e^x-\sum_{j=0}^N\dfrac{x^j}{j!}<\dfrac{3^{N+1}}{(N+1)!}\dfrac{1}{1-\dfrac{3}{N+2}}.\]
For $L=\sum_{k=0}^{20}\dfrac{1}{k!}$, we similarly have $L<e<U=L+\dfrac{1}{20\cdot20!}$. Monotonicity gives a lower bound by the order-$30$ exponential sum at $L$, and an upper bound by that sum at $U$ plus the displayed tail with $N=30$. All quantities are rational, so their comparison is exact. Evaluating gives $15.1542<e^e<15.1543$, hence $\boxed{15.154}$. The following reproduces the bounds without floating-point assumptions:
\begin{verbatim}
from fractions import Fraction as F
from math import factorial
L = sum(F(1,factorial(k)) for k in range(21))
U = L + F(1,20*factorial(20))
def partial(x):
    return sum(x**j/factorial(j) for j in range(31))
tail = F(3**31,factorial(31))/(1-F(3,32))
assert F(151542,10000) < partial(L)
assert partial(U)+tail < F(151543,10000)
\end{verbatim}
